\documentclass[11pt]{article}
\usepackage[margin=1in]{geometry}
\usepackage{amsmath,amssymb,amsthm,booktabs,hyperref}
\newtheorem{lemma}{Lemma}
\newtheorem{proposition}{Proposition}
\newtheorem{theorem}{Theorem}
\newcommand{\ii}{\mathrm i}
\newcommand{\Z}{\mathbb Z}
\title{A linear exact-arithmetic guard for the two-stage complex recursion}
\author{A proposed conditional refinement of the integer-multiplication construction}
\date{October 8, 2026}
\begin{document}
\maketitle

\paragraph{Scope.}
This note improves the precision required for an inherited exact complex network.
It assumes the phase-frame identity, whole-residual recursion and copied
retained-centre schedule of the construction in
\texttt{Swapnil-jain/integer-mult-kappa}, commit
\texttt{f2176bc}, and the interfaces on which that construction depends.
It does not independently establish its Gaussian inverse, layout, prime-selection
or full multiplication theorem. The finite checks accompanying this note check
scalar envelopes and the guard inequalities; they do not replace those assumptions.

\paragraph{Result.}
For the $h=16$ two-stage complex network, $m=256$, the guard may be taken as
\[
                 \boxed{\Delta=128(d+1).}
\]
The same bound also covers the $h=18$ and $h=24$ instances verified below.
In particular the guard exponent is $C_1=1$. The coefficient choice
$b'=128b^3$ and $p=6b'=768b^3$, $b=\lceil\log_2 n\rceil$, meet this guard for
$d=\Theta(b^\epsilon)$ and fixed $\epsilon<1$.
The earlier conservative choice used $p=\Theta(b^{14696})$.
This removes that particular source of an enormous threshold; it is not a
runtime measurement or a proof that the complete algorithm is practical.

\section{The distinction between intermediate and returned values}
Write
\[
 C=\frac{1+\ii}{2}I+\frac{1-\ii}{2}X.
\]
A completed $t$-axis recursive call applies $C^{\otimes t}$, up to address
permutations and inverse/fourth-root wrappers. Its entries belong to
$2^{-t}\Z[\ii]$ and its induced infinity norm is at most $2^t$.
Consequently its returned values consume at most $t$ extra fractional bits
and $t$ magnitude bits, however many arithmetic operations were performed
inside the call. Exact fixed-point representations already contain zeros in
the unused low positions. No truncation, rounding or numerical approximation
is required to exploit this fact.

A depth-first computation has only one currently active descendant at each
level. The extra precision needed inside completed siblings therefore does
not accumulate. For example, concatenating all children only at their
\emph{returned} formats gives a guard recurrence
\[
 D(e)\le E+\frac{s}{m}e+\max_j D(t_j),\qquad
 t_j=w_j\lfloor e/m\rfloor,
\]
where $E$ is a fixed scalar charge, $\sum_jw_j=s$ and
$w_j\le m-h$. This alone yields a linear guard with a large fixed
coefficient. The next argument considerably reduces that coefficient.

\section{A bound for every paused parent prefix}
Let $H_e$ be the normalized Walsh matrix on $e$ bits. For an arbitrary
function $\phi:\mathbb F_2^e\to\Z/4\Z$, put
\[
 F_\phi=H_e\operatorname{diag}(\ii^{\phi})H_e.
\]
The phase-frame invariant of the inherited network places every role in such
a frame. Within each stage, scalar gates act on roles at a common frame. Copied centres are
included: a copy has the same logical value, its phase conversion changes
only its frame, and erasing it discards a temporary rather than an original
scratch role.

\begin{lemma}[Frame ratios]
Every entry of $F_\phi F_\psi^{-1}$ lies in $2^{-e}\Z[\ii]$, and
\[
               \|F_\phi F_\psi^{-1}\|_\infty\le 2^{e/2}.
\]
\end{lemma}
\begin{proof}
The ratio is $H_e\operatorname{diag}(\ii^{\phi-\psi})H_e$.
Expanding its entries gives $2^{-e}$ times a sum of fourth roots of unity
with signs. It is unitary, so every row has Euclidean norm one and has
at most $2^e$ entries. Cauchy--Schwarz proves the infinity-norm bound.
\end{proof}

Refine a scalar prefix to include each addition, multiplication by a fixed
integer, halving, copy and erase operation. Let $M$ be its scalar matrix
in logical coordinates. This matrix can be rectangular when a temporary is
created or removed. The block from initial role $i$ to current role $j$ is
\[
                 M_{ji}F_{\mathrm{current},j}F_{\mathrm{initial},i}^{-1}.
\]
Let $B$ bound the denominator exponent of every entry of every such $M$,
and let $G$ bound all its absolute row sums. The lemma shows that this
prefix increases the fractional exponent by at most $e+B$ and the
base-two magnitude bound by at most $e/2+\log_2G$.
This statement allows different initial frames on different roles and
arbitrary initial scratch values.

A partial residual edge, after any completed whole child blocks, still has
an arbitrary Walsh-diagonal phase frame. The coordinate changes used to
implement an edge are address permutations, and so do not increase either
numerical bound. Thus the same estimate holds immediately before every
child call, not merely between scalar gates.

\section{Finite scalar envelope}
The program \texttt{semantic\_guard.py} builds the actual producer and
compiled role program. It assigns every scalar input row norm one and
propagates nonnegative absolute row-sum bounds through every addition.
Inverse words reverse the compiled operations and use their inverse order;
ignoring signs is conservative. Integer numerators are charged before a
halving, so intermediate products are included. Denominator exponents
propagate by taking a maximum and adding one at a factor $1/2$.

The retained scatter requires particular care. Let
$T=\sum_t x_t$ and $E_i=\sum_{t\not\ni i}x_t$, retaining $E_i$ for
$i<h-1$. Its value for target triple $S$ is
\[
 \begin{cases}
 T-\frac12\sum_{i\in S}E_i,&h-1\notin S,\\[2pt]
 \frac12\bigl((5-h)T+\sum_{i<h-1,\ i\notin S}E_i\bigr),&h-1\in S.
 \end{cases}
\]
The second formula follows from $\sum_{i=0}^{h-1}E_i=(h-3)T$.
Some of its coefficients have magnitude larger than one; the verifier
includes them explicitly.

Stage-one invocations have disjoint banks, and the same holds for stage two.
Their uniform envelopes therefore compose by multiplication, rather than
by raising them to the number of invocations. The inherited endpoint
correction copies output $A$, applies the inverse rank-one phase child
$P^{-1}$ to its copy, adds it to output $B$, and erases the copy.
Indeed its logical expressions are
\[
 A=-Fy,\qquad B=FP^{-2}x+Ey,\qquad E=FP^{-1},
\]
so $B+P^{-1}A=FP^{-2}x$. The copy and erase add no arithmetic growth,
and the one addition costs at most a factor two in the scalar envelope.
Its phase child is already included in the recursive rank histogram.
The endpoint addition is handled separately from the common-frame assertion:
multiplying the copied row by $P^{-1}$ preserves the family of Walsh-diagonal
phase operators. Adding it to the other row adds their scalar row-sum bounds,
without enlarging the common $2^{-e}\Z[\ii]$ denominator grid. No alignment
of the two endpoint frames is needed for this separate bound.
Signs and fourth-root corrections preserve these bounds.

The exact envelope results are
\[
\begin{array}{c|r|r|r|r}
 h&G_{\mathrm{forward}}&G_{\mathrm{inverse}}&
       G=2G_{\mathrm{forward}}G_{\mathrm{inverse}}&B\\\hline
16&135386&256057&69333066004&2\\
18&250444&476800&238823398400&2\\
24&1006577&1925603&3876535381862&2
\end{array}
\]
All satisfy $G<2^{48}$ and $B\le8$. We deliberately use the looser
common bound
\[
               \text{paused parent guard}\le 2e+64.       \tag{1}
\]
This includes the scalar intermediates in the verifier, rather than just
the final scalar outputs.

\section{Integer widths and the linear induction}
The integer-width interface supplies children of width
$w\lfloor e/m\rfloor$, with $w\le m-h$.
For an input width not divisible by $m$, first evaluate its fewer than $m$
remainder axes individually and apply the network to the other
$m\lfloor e/m\rfloor$ axes. A completed individual kernel costs one
semantic fractional bit and at most one magnitude bit; its own
intermediates use the usual eight-bit-per-axis allowance.
These remainder operations introduce only additional phase frames in the
prefix argument. They do not introduce new axes. For fewer than $m$ axes,
or at the original stopping threshold, evaluate all kernels individually.

For the divisible internal call, (1) and the one-active-child observation give
\[
 D(e)\le 2e+64+\max_j D(t_j),\qquad
 t_j\le\frac{h-1}{h}e,\quad e\ge h^2.
\]
Choose
\[
 C\ge8,\qquad \frac{C}{h}\ge 2+\frac{64}{h^2}.
\]
Then induction proves $D(e)\le Ce$: the second inequality makes
$2e+64+C(h-1)e/h\le Ce$, and the first covers leaves. A remainder of
$r$ individually evaluated axes followed by a block of $e-r$ axes
costs at most $\max\{8r,r+C(e-r)\}\le Ce$, so the same bound covers
all integer widths. Verified integer choices are
\[
       (h,C)=(16,36),\quad(18,40),\quad(24,51).
\]
This induction is independent of the stopping exponent $\beta$.

\begin{theorem}[Whole-layer guard]
Under the inherited exact interfaces and the scalar/phase invariants above,
$\Delta=128(d+1)$ suffices for one normalized butterfly layer in each
of the three verified instances.
\end{theorem}
\begin{proof}
Top-level active-axis pieces are disjoint and their total width is at most
$d$. Completed earlier pieces apply the exact tensor phase operator, so
collectively they consume at most $d$ semantic guard bits.
During the currently active piece its extra guard is at most $Cd$.
Individually handled preprocessing and tail axes contribute at most $8d$;
this conservatively charges their intermediates as well as their returns.
The outer phase factors preserve magnitude and denominator. The conversion
$\mathfrak b^D$, with $\mathfrak b=(1-\ii)/2$, uses at most $d+9$ further guard bits,
including its binary shift and its possible final factor $\mathfrak b$.
Thus $(C+10)d+9$ suffices. Since $C\le51$, this is below $128(d+1)$.
Padding, splitting, copying and changes of address basis do not perform
arithmetic on coefficient values. The usual fixed-width exact format with
$p+\Delta$ fractional bits and $\Delta+3$ integer/sign bits therefore
represents every intermediate value without rounding or overflow.
\end{proof}

\paragraph{Cubic precision with a finite numerical margin.}
Use the longer digit width $b'=128b^3$ and precision $p=6b'=768b^3$;
in the notation of the inherited assembly this is $x=2$.
Set $d=\lfloor b^\epsilon\rfloor$ and $y=\epsilon$, with $0<\epsilon<1$.
For every integer $b\ge1$, the guard itself satisfies
\[
 \Delta=128(d+1)\le256b\le768b^3=p.
\]
This inequality is a numerical guard cutoff only: the complete layout and
prime-selection hypotheses still have to hold.

The inherited Gaussian choices are
\[
 \alpha^2=\left\lfloor\frac{p}{48d}\right\rfloor
          =\left\lfloor\frac{16b^3}{d}\right\rfloor,
 \qquad \eta=\frac1{4dL^\epsilon}.
\]
Once the inherited prime-selection step supplies the constant-factor band
$\theta_i\ge\eta/2$, use $d\le b^\epsilon$ and $L\le b$.
Since $b^3/d\ge1$, the floor costs at most one unit and
\[
 \alpha^2\ge\frac{15b^3}{d},\qquad
 \alpha^2\theta_i\ge
 \frac{15b^3}{8d^2L^\epsilon}
 \ge\frac{15}{8}b^{3-3\epsilon}\ge\frac{15}{8}>1.
\]
Thus no eventual domination of a factor $1/32$ is needed for this particular
Gaussian inequality. Moreover
$\gamma=2d\alpha^2\le32b^3=b'/4$, as required by the inherited final
error estimate. These bounds do not establish the remaining hypotheses of
that estimate.

The guard condition is now $\epsilon<3$, and the inherited segmented-inverse
precision condition is $3\epsilon<3$, both automatic for fixed
$\epsilon<1$. A factor $p^\delta$ contributes exponent $3\delta$ in powers
of $b$. Choices of $\epsilon,\delta,\beta$ must still satisfy the final
saving and all other assembly inequalities.

\paragraph{Inherited prime intervals and remaining thresholds.}
The prime-interval argument is the one in
\href{https://github.com/Swapnil-jain/integer-mult-kappa/blob/f2176bc/notes/stack-notes.tex}
{\texttt{notes/stack-notes.tex}, Section 1}; it is not a new prime-gap result
proved here. It uses the classical prime number theorem with an error term
$O(X\exp(-c\sqrt{\log X}))$, for a positive constant $c$.
Here $\log X=\Theta(L/d)$ and the relative interval length
$\eta=1/(4dL^\epsilon)$ is a fixed inverse power of $L$, hence a fixed
inverse power of $\log X$ for each fixed $\epsilon<1$.
The interval main term $\eta X/\log X$ eventually dominates that error
and supplies more than $d$ primes in each required interval. The constant-factor band
used above is part of that inherited eventual selection argument.
This note supplies no explicit common cutoff for the prime intervals,
record lengths, axis widths, or remaining asymptotic estimates. In
particular, a small numerical guard cutoff does not make the whole
multiplication algorithm practical or yield a benchmark.

\paragraph{Reproduction and attribution.}
From the accompanying \texttt{research} directory, run
\texttt{python3 independent/guard-improvement/semantic\_guard.py 16 18 24}
and the adjacent unit tests. The phase-frame construction, scalar
producer, copied-centre schedule and endpoint transfer are inherited;
this note's refinement is to charge returned semantic values and use
phase-prefix cancellation to prove a linear exact-arithmetic guard.
\end{document}
